\section{Analytic framework and the credited dependency chain}
\label{sec:analytic}

The analytic corpus supplies exact identities and explicit approximation
bounds. Its recorded Judge receipts distinguish successful module rows from
the status of an entire batch. This distinction matters for batches 6, 7, 8,
9, 11 and 26: each globally failed batch retains only the whole-module PASS
rows identified in the inventory. No failed or uninvoked module is credited.
The standard axioms recorded for credited declarations are subsets of
\texttt{propext}, \texttt{Classical.choice} and \texttt{Quot.sound}.
Appendix~\ref{app:analytic-inventory} provides a module table; the accompanying
JSON and CSV retain every extracted axiom-audit row.\ev{A-PARTIAL-BATCHES}
\ev{A-INVENTORY-COUNTS}

\subsection{The genuine coefficient and its continuous projection}

Throughout this section, $\Lambda$ is the genuine von Mangoldt function:
$\Lambda(0)=\Lambda(1)=0$, $\Lambda(p^k)=\log p$ for a prime $p$ and
an integer $k\geq1$, and it vanishes otherwise. For $a>0$, set
\begin{equation}
 T_a(\theta)=\sum_{n\geq0}\Lambda(n)e^{-an}e^{in\theta},\qquad
 C_N=\sum_{n=0}^{N}\Lambda(n)\Lambda(N-n).
 \label{eq:analytic-definitions}
\end{equation}
The series is absolutely convergent. The compiled continuous projection
identity, for every natural $N$, is
\begin{equation}
 C_N=\frac{e^{aN}}{2\pi}\int_0^{2\pi}
          T_a(\theta)^2e^{-iN\theta}\,d\theta.
 \label{eq:circle-identity}
\end{equation}
This is \nolinkurl{thermalContinuousProjection_eq_vonMangoldt_sum}
in the module \nolinkurl{ThermalProjectionIdentity22}.
The source uses the full interval $[0,2\pi]$. Prime powers remain in $C_N$;
the equality alone is not a lower bound for a prime-pair sum. The square
$T_a^2$ extracts additive pairs. Replacing it by $|T_a|^2$ extracts
differences and changes the coefficient being studied.\ev{A-CIRCLE}
\ev{A-OBSTRUCTION-SIGNED-GAP}

There is also a compiled elementary truncation envelope. With $r=e^{-a}$,
let
\begin{equation}
 U(a)=\frac{r}{(1-r)^2},\qquad
 B(a,M)=r^{M+1}\left(\frac{M+1}{1-r}+U(a)\right).
\end{equation}
If $T_{a,M}$ is the sum through $M$, the corresponding normalized
projection differs from the full projection by at most
$2e^{aN}U(a)B(a,M)$ in complex norm. The bound follows from actual
von Mangoldt terms and a constructed summable majorant, rather than from
an assumed coefficient bound. Its module is a retained PASS row in failed
batch 11.\ev{A-CIRCLE}

\subsection{Finite projection and canonical logarithm quantization}

For natural $N\leq M$ and $K>\max(N,2M-N)$, the discrete projector satisfies
\begin{equation}
 C_N=\frac{e^{aN}}{K}\sum_{j=0}^{K-1}
 T_{a,M}(2\pi j/K)^2 e^{-2\pi iNj/K}.
 \label{eq:discrete-projector}
\end{equation}
The condition excludes all unwanted exponents congruent to $N$ modulo $K$.
This is a finite identity for every real $a$, including $a=0$; it does not
assert convergence of the infinite series at $a=0$.
The declaration \nolinkurl{sampled_trueCircle_eq_coefficient}
in \nolinkurl{DiscreteThermalProjection22}
records the concrete angular grid.\ev{A-DISCRETE}

The quantized construction uses the fixed scale $S=2^{58}$. Define
\begin{equation}
 L_{32}(z)=2\sum_{j=0}^{31}\frac{z^{2j+1}}{2j+1},\qquad
 \rho_{32}=\frac{9}{4\cdot65\cdot3^{65}}.
\end{equation}
For $2\leq p\leq10^8$, write $k=\lfloor\log_2p\rfloor$ and
$z=(p-2^k)/(p+2^k)$. The rational lower endpoint is
$b_p=kL_{32}(1/3)+L_{32}(z)$, with interval width
$(k+1)\rho_{32}$. The constructor rounds the midpoint multiplied by $S$
to the nearest integer, resolving exact ties toward the even integer,
then clamps to $[0,32S]$. The proved error is at most $1/S$;
precision is a conclusion of this fixed constructor.\ev{A-LOG}

Applying the constructor to the base prime of each prime power defines
integers $A_n$ and $q_n=A_n/S$. Both $q_n$ and $\Lambda(n)$ lie in $[0,32]$
on $0\leq n\leq10^8$, and $|q_n-\Lambda(n)|\leq1/S$ there. Thus, for
natural $N\leq10^8$, the canonical integer coefficient
$I_N=\sum_{n=0}^N A_nA_{N-n}$ satisfies
\begin{equation}
 \left|C_N-\frac{I_N}{S^2}\right|
 \leq E_{\log}(N):=(N+1)\left(\frac{64}{S}+\frac1{S^2}\right).
 \label{eq:log-envelope}
\end{equation}
The fixed integer guard proves $2E_{\log}(10^8)\leq10^{-6}$.
These statements are in \nolinkurl{RationalLogQuantization22} and
\nolinkurl{QuantizedLambdaEnvelope22}; no real-logarithm precision
assumption is inserted.\ev{A-LOG}

\subsection{Finite-field extraction and the five concrete roots}

The generic finite-field module derives character orthogonality from an
actual primitive root, then instantiates the coefficient projector for the
canonical integer sequence. The concrete root module proves the primality
of the five moduli and exact order $K=2^{27}$ of the roots
$\omega=g^{(p-1)/K}$ in their respective fields. Here $g$ denotes the
specified bank base; the needed assertion is the order of $\omega$.
\ev{A-FIELD}

\begin{center}
\begin{tabular}{rr}
\toprule Certified prime $p$ & Bank base $g$\\\midrule
2013265921 & 31\\
2281701377 & 3\\
3221225473 & 5\\
3489660929 & 3\\
3892314113 & 3\\
\bottomrule
\end{tabular}
\end{center}

At $N=M=10^8$, each such field projector returns $I_N$ modulo $p$.
The compiled results are in \nolinkurl{FiniteFieldProjection22},
\nolinkurl{ConcreteNTTRoots22} and
\nolinkurl{ConcreteNTTA32Projection22}. A theorem about this finite
projector does not by itself refine the native transform or CRT program.
The saved numerical computation is reported separately in
Section~\ref{sec:numerical}.\ev{A-FIELD}

\subsection{The actual Gamma inversion dependency chain}

The inversion chain is a sequence of credited source modules with explicit
imports, rather than an assumed Gamma inversion formula. The closed-strip
prerequisite bounds the actual Gamma function by
$2\exp(-\pi|\operatorname{Im}s|/4)$ for
$1\leq\operatorname{Re}s\leq2$. The positive-real inverse is first proved;
the local complex module constructs its branch and integrable majorants;
the holomorphy module then extends the identity to the right half-plane.
\ev{A-GAMMA-CHAIN}

\begin{center}\small
\begin{tabular}{@{}p{0.48\textwidth}p{0.12\textwidth}p{0.29\textwidth}@{}}
\toprule Module & Batch; rows & Retained scope\\\midrule
\nolinkurl{GammaPrerequisites22} & 2; 23 & Actual rotation and strip bound\\
\nolinkurl{ThermalGammaMellinInverse22} & 20; 11 & Positive-real inversion\\
\nolinkurl{ComplexGammaMellinLocal22} & 26; 22 & Local bounds and real agreement; partial PASS\\
\nolinkurl{ComplexGammaMellinHolomorphy22} & 27; 15 & Complex inversion for $\operatorname{Re}w>0$\\
\nolinkurl{ComplexGammaMellinLambda22} & 35; 30 & Actual infinite $\Lambda$ interchange\\
\nolinkurl{ComplexGammaMellinTail22} & 30; 22 & Unweighted Gamma tail\\
\nolinkurl{ComplexGammaMellinExpTailBridge22} & 32; 2 & Exponential inversion error\\
\bottomrule
\end{tabular}
\end{center}

The first five rows follow the import chain in their displayed order.
Batch 26 failed globally, but its entire local inversion module has the
retained PASS row. The unweighted tail module imports the local module;
the exponential-tail bridge imports both holomorphy and tail. The
$\Lambda$ interchange imports holomorphy. It does not require the
uncompiled weighted-$\Lambda$ tail module. Counts in this table are saved
axiom-audit declaration rows, including definitions.\ev{A-GAMMA-CHAIN}

\subsection{The genuine infinite sum--integral interchange}

For real $t$ and complex $w$ with $\operatorname{Re}w>0$, set
\begin{equation}
 D(t)=\sum_{n\geq2}\Lambda(n)n^{-2-it},\qquad
 P(w)=\sum_{n\geq0}\Lambda(n)e^{-nw}.
\end{equation}
The source defines these actual arithmetic series. It proves continuity
of $D$, the bound $|D(t)|\leq6$, and summability of the integrals of norms
needed for the infinite interchange. In particular, the coefficient
majorant is bounded by $2n^{-3/2}$ and its total mass by 6. The compiled
conclusion is
\begin{equation}
 P(w)=\frac1{2\pi}\int_{\R}D(t)\Gamma(2+it)w^{-2-it}\,dt,
 \qquad \operatorname{Re}w>0,
 \label{eq:true-mellin}
\end{equation}
with principal complex powers. The declaration
\nolinkurl{lambdaMellinThermal_eq_integral}
in \nolinkurl{ComplexGammaMellinLambda22}
is therefore an infinite arithmetic interchange result. This module alone
does not identify $D$ with a zeta logarithmic derivative or supply a
weighted-$\Lambda$ tail bound uniform on the coefficient contour.
\ev{A-MELLIN}

\subsection{Angular integration by parts and the remaining boundary}

For $a>0$, natural $N\geq1$ and $q\in\C$, the compiled angular module
uses
\begin{equation}
 J_{a,N}(q)=\frac1{2\pi}\int_{-\pi}^{\pi}
                  (a-i\theta)^{-q}e^{-iN\theta}\,d\theta.
\end{equation}
Its exact recurrence is
\begin{align}
 J_{a,N}(q)&=B_{a,N}(q)+\frac qN J_{a,N}(q+1),\\
 B_{a,N}(q)&=\frac{i(-1)^N}{2\pi N}
           \bigl((a-i\pi)^{-q}-(a+i\pi)^{-q}\bigr).
 \label{eq:angular-boundary}
\end{align}
The boundary term at $q=1$ is proved nonzero. Periodicity of the complete
arithmetic trace therefore cannot be substituted for periodicity of an
individual principal-branch Mellin kernel. This gives a precise boundary
obstruction to dropping that term; it does not rule out estimates that
retain it.\ev{A-IBP}\ev{A-OBSTRUCTION-BOUNDARY}

The unweighted Gamma-tail and exponential-inversion-error theorems are
compiled. The repaired weighted-$\Lambda$ tail, uniform circle geometry,
coefficient truncation envelope, scalar connectors and coefficient-to-scalar
bridge are retained source only. Failed batch 36 adds no credit: the
weighted-tail module failed, and the geometry module was not invoked.
Consequently the complete Mellin-to-coefficient truncation chain remains
uncompiled. A norm enclosure, even if supplied by that chain, would still
leave the signed estimate and prime-power corrections needed for a
prime-pair conclusion.\ev{A-TRUNCATION-BOUNDARY}
\ev{A-OBSTRUCTION-SIGNED-GAP}
